\clearpage
\begin{explanationbox}
\textbf{先看一组实数：$AB=4,AC=6$。}
\begin{center}
\begin{tikzpicture}[font=\small]
\draw[black!30,rounded corners=2pt] (0,0) rectangle (14.4,8.2);
\node[anchor=west,font=\small\bfseries] at (0.5,7.72) {$AB<AC$：外分点在 $B$ 的外侧，顺序为 $D,B,C$};
% MODEL 4 6 4
\begin{scope}[shift={(7.5,3)},scale=0.8]
\coordinate (A) at (-0.5,3.9686269666);
\coordinate (B) at (0,0);
\coordinate (C) at (4,0);
\coordinate (D) at (-8,0);
\coordinate (X) at (-1.34375,4.7127445228);
\draw[black!45] (-8.6,0)--(4.6,0);
\draw[blue,thick] (A)--(B);
\draw[teal!80!black,thick] (A)--(C);
\draw[black,thick] (B)--(C);
\draw[purple,very thick] (A)--(D);
\fill (A) circle (1.5pt) node[above=4pt] {$A$};
\fill (B) circle (1.5pt) node[below left=4pt] {$B$};
\fill (C) circle (1.5pt) node[below right=4pt] {$C$};
\fill (D) circle (1.5pt) node[below right=4pt] {$D$};
\draw[black!45,dashed] (A)--(X);
\path (A)--(B) node[midway,above left=3pt] {$AB=4$};
\path (A)--(C) node[midway,above right=3pt] {$AC=6$};
\draw[purple,thick] (-1.0303300859,3.6880026626) arc[start angle=-152.1144331639,delta angle=69.2951889454,radius=0.6];
\draw[purple,thick] (-1.08125,4.4812412831) arc[start angle=138.5903778907,delta angle=69.2951889454,radius=0.775];
\end{scope}
\draw[blue,|-|] ($(D)+(0,-1.1)$)--($(B)+(0,-1.1)$) node[midway,below=3pt] {$BD=8$};
\draw[teal!80!black,|-|] ($(D)+(0,-2.05)$)--($(C)+(0,-2.05)$) node[midway,below=3pt] {$DC=12$};

\end{tikzpicture}
\end{center}

因为 $AB/AC=2/3<1$，必须有 $BD<DC$。
外点只能在 $B$ 的外侧：$D,B,C$ 依次排列，$DC$ 比 $BD$ 多出整段 $BC$。
图中 $8/12=4/6$，且 $12-8=4$。
虚线延长一条邻边，紫色两段圆弧标出所平分的外角。

\textbf{换一组邻边：$AB=5,AC=3$。}
\begin{center}
\begin{tikzpicture}[font=\small]
\draw[black!30,rounded corners=2pt] (0,0) rectangle (14.4,8.2);
\node[anchor=west,font=\small\bfseries] at (0.5,7.72) {$AB>AC$：外分点在 $C$ 的外侧，顺序为 $B,C,D$};
% MODEL 5 3 4
\begin{scope}[shift={(2,3)},scale=0.8]
\coordinate (A) at (4,3);
\coordinate (B) at (0,0);
\coordinate (C) at (4,0);
\coordinate (D) at (10,0);
\coordinate (X) at (4.9,3.675);
\draw[black!45] (-0.6,0)--(10.6,0);
\draw[blue,thick] (A)--(B);
\draw[teal!80!black,thick] (A)--(C);
\draw[black,thick] (B)--(C);
\draw[purple,very thick] (A)--(D);
\fill (A) circle (1.5pt) node[above=4pt] {$A$};
\fill (B) circle (1.5pt) node[below left=4pt] {$B$};
\fill (C) circle (1.5pt) node[below right=4pt] {$C$};
\fill (D) circle (1.5pt) node[below right=4pt] {$D$};
\draw[black!45,dashed] (A)--(X);
\path (A)--(B) node[midway,above left=3pt] {$AB=5$};
\path (A)--(C) node[midway,below right=4pt] {$AC=3$};
\draw[purple,thick] (4.5366563146,2.7316718427) arc[start angle=-26.5650511771,delta angle=63.4349488229,radius=0.6];
\draw[purple,thick] (4,2.225) arc[start angle=-90,delta angle=63.4349488229,radius=0.775];
\end{scope}
\draw[blue,|-|] ($(D)+(0,-1.1)$)--($(B)+(0,-1.1)$) node[midway,below=3pt] {$BD=10$};
\draw[teal!80!black,|-|] ($(D)+(0,-2.05)$)--($(C)+(0,-2.05)$) node[midway,below=3pt] {$DC=6$};

\end{tikzpicture}
\end{center}

现在 $BD/DC=5/3>1$，$D$ 在 $C$ 的外侧，$BD$ 比 $DC$ 多出 $BC$。
两种图的长度单位一致；它们替代了“延长线到底是哪一边”的文字猜测。

\textbf{为什么两条邻边接近时，交点可能很远？}
令 $r=AB/AC$，并用 $BC$ 作长度单位。
此时 $B$ 的坐标为 $0$、$C$ 为 $1$，外分点坐标为 $t=r/(r-1)$。
\begin{center}
\begin{tikzpicture}[font=\small]
\draw[black!30,rounded corners=2pt] (0,0) rectangle (14.4,8.2);
\node[anchor=west,font=\small\bfseries] at (0.5,7.72) {邻边越接近，外分点越远：横轴 $r=AB/AC$，纵轴 $t=x_D/BC$};

\begin{axis}[width=11.2cm,height=6.1cm,at={(2.05cm,1.25cm)},anchor=south west,
axis lines=left,xlabel={$r$},ylabel={$t$},xmin=.5,xmax=1.5,ymin=-12,ymax=12,
xtick={.5,.75,1,1.25,1.5},ytick={-10,-5,0,5,10},
tick label style={font=\small},label style={font=\small},
legend style={font=\small,draw=none,fill=white,at={(.03,.97)},anchor=north west}]
\addplot[draw=none,fill=black!8,forget plot] coordinates {(.5,0)(1.5,0)(1.5,1)(.5,1)}\closedcycle;
\addplot[purple,very thick,domain=.5:.95,samples=100] {x/(x-1)};\addlegendentry{$t=r/(r-1)$}
\addplot[purple,very thick,domain=1.05:1.5,samples=100,forget plot] {x/(x-1)};
\draw[black!60,densely dashed] (axis cs:1,-12)--(axis cs:1,12);
\node[anchor=west] at (axis cs:1.13,9) {$r\to1,\ |t|\to\infty$};
\node[anchor=west] at (axis cs:.53,3) {$0<t<1$ 是 $BC$ 内部};
\end{axis}

\end{tikzpicture}
\end{center}

$r=0.8$ 时 $t=-4$，$r=1.2$ 时 $t=6$：交点都在 $[0,1]$ 外。
当 $r\to1$ 时，分母趋于零，$|t|\to\infty$。
图中的竖虚线 $r=1$ 表示公式无定义；真正等腰时是两条直线平行，不能画成一个很远但有限的交点。
\end{explanationbox}
